\documentclass[11pt]{article}
\usepackage[margin=1in]{geometry}
\usepackage{amsmath,amssymb,amsthm}
\usepackage{graphicx}
\usepackage[hidelinks]{hyperref}

\newtheorem{theorem}{Theorem}
\newtheorem{remark}{Remark}

\newcommand{\R}{\mathbb R}
\newcommand{\calL}{\mathcal L}

\title{A Formal Path Dilation for Every Finite-Horizon Evolution Family}
\author{Run fa\_banach\_001, agent\_lane\_19}
\date{June 20, 2026}

\begin{document}
\maketitle

\noindent\textbf{Status.} Partial result / clarification, likely valid. The
source asks which \(C_0\)-semigroups and \(C_0\)-evolution families admit an
approximate invertible dilation. Read literally, without requiring the dilation
space to preserve Hilbertianity, \(2\)-smoothness, UMD, or related geometry,
the answer is: all finite-horizon \(C_0\)-evolution families admit an exact
invertible dilation. The construction is too large for the stochastic maximal
inequality theorem in the source, so this is not a solution of the intended
geometry-preserving problem.

\section*{Source Question}

The source is Jan van Neerven and Mark Veraar, \emph{Maximal estimates for
stochastic convolutions in \(2\)-smooth Banach spaces and applications to
stochastic evolution equations}, arXiv:2006.08325. In Section 6, the authors
ask:
\begin{quote}
``Which \(C_0\)-semigroups and \(C_0\)-evolution families admit an
\((\)approximate\()\) invertible dilation?''
\end{quote}
They explain just above and below the question that the useful version for
maximal inequalities requires the dilation spaces to retain the relevant
geometry, such as Hilbertianity, \(2\)-smoothness, or UMD.

\begin{center}
\includegraphics[width=\linewidth]{figures/open_question_crop.png}
\end{center}

\section*{Definition from the Source}

Let \((S(t,s))_{0\leq s\leq t\leq T}\) be a \(C_0\)-evolution family on a
Banach space \(X\). In the source terminology, it admits an invertible dilation
on a Banach space \(Y\) if there are strongly continuous functions
\[
  J:[0,T]\to\calL(X,Y),
  \qquad
  Q:[0,T]\to\calL(Y,X)
\]
such that
\[
  S(t,s)=Q(t)J(s),\qquad 0\leq s\leq t\leq T.
\]

\section*{Idea of the Proof}

The factorization can be forced by remembering the future orbit of \(x\) from
the starting time \(s\). Put this orbit into the path space \(C([0,T];X)\), and
make it constant before time \(s\) so that the path is continuous at the
switching point. Evaluation at \(t\) then recovers \(S(t,s)x\). Strong
continuity of the maps \(J(s)\) follows from the uniform strong continuity of
the evolution family on the compact time triangle.

\begin{theorem}[Formal path dilation]\label{thm:path}
Let \(X\) be a Banach space and let
\((S(t,s))_{0\leq s\leq t\leq T}\) be a \(C_0\)-evolution family on \(X\).
Then \((S(t,s))\) admits an exact invertible dilation, in the literal sense of
the source definition, on \(Y=C([0,T];X)\).
\end{theorem}

\begin{proof}
Let \(Y=C([0,T];X)\) with the supremum norm. For \(t\in[0,T]\), define
\[
  Q(t)f=f(t),\qquad f\in Y.
\]
Then \(\|Q(t)\|\leq 1\), and \(t\mapsto Q(t)\) is strongly continuous because
each \(f\in C([0,T];X)\) is uniformly continuous.

For \(s\in[0,T]\) and \(x\in X\), define \(J(s)x\in Y\) by
\[
  (J(s)x)(r)=
  \begin{cases}
    x, & 0\leq r\leq s,\\
    S(r,s)x, & s\leq r\leq T.
  \end{cases}
\]
This is a continuous \(X\)-valued path, since \(S(s,s)=I\). The family
\(\{S(t,s):0\leq s\leq t\leq T\}\) is pointwise bounded by strong continuity
on the compact triangle, hence uniformly bounded by the uniform boundedness
principle. Therefore
\[
  \|J(s)x\|_{\infty}\leq
  \max\left\{1,\sup_{0\leq a\leq b\leq T}\|S(b,a)\|\right\}\|x\|,
\]
so \(J(s)\in\calL(X,Y)\) with a bound uniform in \(s\).

It remains to check strong continuity of \(s\mapsto J(s)\). Fix \(x\in X\) and
let \(s_n\to s\). The map
\[
  (r,a)\mapsto S(r,a)x
\]
is continuous on the compact triangle \(0\leq a\leq r\leq T\), hence uniformly
continuous there. If \(r\geq \max\{s_n,s\}\), uniform continuity gives
\[
  \|S(r,s_n)x-S(r,s)x\|\to 0
\]
uniformly in \(r\). If \(r\leq \min\{s_n,s\}\), the two paths are both equal
to \(x\). On the remaining interval between \(s_n\) and \(s\), one path is
equal to \(x\), while the other is \(S(r,a)x\) with \(|r-a|\leq |s_n-s|\);
uniform continuity at the diagonal, together with \(S(a,a)=I\), gives
\[
  \sup_{|r-a|\leq |s_n-s|}\|S(r,a)x-x\|\to 0.
\]
Thus \(\|J(s_n)x-J(s)x\|_{\infty}\to0\), proving strong continuity.

Finally, for \(0\leq s\leq t\leq T\),
\[
  Q(t)J(s)x=(J(s)x)(t)=S(t,s)x.
\]
This is the desired exact dilation.
\end{proof}

\begin{remark}[Why this does not settle the geometric problem]
The construction generally destroys the geometry needed in the source's
dilation theorem. If \(X\neq\{0\}\), choose \(x_0\in X\) with \(\|x_0\|=1\).
The map \(f\mapsto (r\mapsto f(r)x_0)\) embeds \(C([0,T])\) isometrically into
\(C([0,T];X)\). Hence \(C([0,T];X)\) is non-reflexive. In particular it cannot
be \(2\)-smooth or UMD, since those classes are reflexive. It is also not a
Hilbert space except in the trivial zero-dimensional case.

Thus Theorem~\ref{thm:path} answers the literal non-geometric version of the
dilation question, but the source's useful question remains: when can such a
factorization be arranged inside a class preserving Hilbertian, \(2\)-smooth,
UMD, or other stochastic-integration geometry?
\end{remark}

\section*{Verification and Novelty Check}

Local indexes were searched for arXiv:2006.08325 and for the phrases
\emph{invertible dilation}, \emph{approximate invertible dilation},
\emph{evolution family}, and \emph{stochastic convolution}. No existing packet
recorded this formal path-space factorization. Web searches on June 20, 2026
for combinations of \emph{invertible dilation}, \emph{\(C([0,T];X)\)}, and
\emph{evolution family} did not find a later answer to the source question.

\section*{Human Review Recommendation}

Review should treat this as a clarification packet, not as a full stochastic
maximal-inequality solution. The proof of the factorization is elementary; the
main review point is whether the run should record the source question as
implicitly geometry-preserving. If so, this packet should remain in
\texttt{partial/}: it identifies the formal obstruction and sharpens the
remaining problem.

\begin{thebibliography}{9}
\bibitem{NV}
J. van Neerven and M. Veraar,
\emph{Maximal estimates for stochastic convolutions in \(2\)-smooth Banach
spaces and applications to stochastic evolution equations},
arXiv:2006.08325, 2020.
\end{thebibliography}

\end{document}
